\documentclass{article}
\usepackage[utf8]{inputenc}
\usepackage[T1]{fontenc}
\usepackage[margin=2cm]{geometry}
\usepackage{hyperref}
\sloppy
\begin{document}
\section*{Historical source}
Encyclopaedia Britannica scan, source row 1327. Parse 3.5 extracted the text. Super 3.5 VL corrected it from an enlarged crop of each block. It can still contain recognition errors.
Page headers, footers, and page numbers are removed.

\subsection*{Untitled continuation (E0109)}
\textit{Continues an entry from an earlier page.}

Hier.
Hier.
Hucft.

der begins to fill; it rises gradually to the surface, and gradually encreases its height, smoking amazingly, and slinging up great stones. After attaining its greatest height it gradually sinks till it totally disappears. Boiling jets d'eaux and boiling springs are frequent in most parts of the island. In many parts they are applied to the culinary uses of the natives. The most capital is that which is called Geyer or Geyser, in a plain rising into small hills, and in the midst of an amphitheatre, bounded by the most magnificent and various-shaped icy mountains; among which the three-headed Hecla fora pre-eminent. See Iceland, N° 4.—The fires are not confined to the land; they rise in the very sea, and form fountains amidst the waves. Their distance from the land is unknown; but the new volcanic isle, twelve miles off the point of Reickenes, emitting fire and smoke, proves that the subterraneous fires and waters extend to that space; for those awful effects arise from the united fury of the two elements.

\subsection*{HUESCA (E0110)}
HUESCA, an ancient and considerable town of
Spain, in the kingdom of Arragon, with a bishop's
fee and an university. It is seated on the Iffuela, in
a soil producing excellent wine, in W. Long. o. 13.
N. Lat. 40. 2.

\subsection*{HUESCAR (E0111)}
HUESCAR, or GUESCAR, a town of Spain, in the
kingdom of Granada, seated on a plain in W. Long.
1. 45. N. Lat. 37. 32.

\subsection*{HUESNE (E0112)}
HUESNE, or HVENA, a small island in the Baltic Sea, in the Sound, where was the famous observatory of Tycho Brahe. E. Long. 13. 5. N. Lat. 55. 54.

HUET (Peter Daniel), a very learned French writer, born at Caen in Normandy, on the 8th of February 1630. He discovered, from his infancy, a great inclination to the study of polite literature and the sciences, and at first applied himself to the law; but Des Cartes's principles, and Bochart's sacred geography, made him change his studies for those of philosophy, mathematics, the languages, and antiquities. His admiration for Bochart made him desirous of knowing him. He contracted a very strict friendship with him, and accompanied that learned man to Sweden. Here Christina would have engaged him in her service; but he, sensible of her inconstant temper, returned to France. All he brought with him was a copy of a MS. of Origen, which he transcribed at Stockholm. He refused several offers from Christina after she abdicated and went to Rome, and from Gustavus her successor. In 1670, Mr. Boileau being appointed by the king preceptor to the dauphin, his majesty chose Mr. Huet for his colleague, with the title of sub-preceptor to the prince. It was he that formed the plan of the commentaries in usum Delphini, and directed the execution. His sentiments of piety determined him to enter into holy orders, which he did at the age of 46. Soon after this, he was preferred by the king to the abbey of Aunay; and in 1685 was nominated to the bishopric of Soissons, which he exchanged for the see of Avranches. After governing that diocese ten years, he resigned, and was made abbot of Fontenay near Caen. His love to his native place determined him to fix there. But lawsuits coming upon him, he retired to Paris, and lodged among the Jesuits in the Maison Professe, whom he

had made heirs to his library. A severe distemper
weakened his body extremely, but not the vivacity of
his genius: he wrote his own life in a very elegant
style; and died in 1721, aged 91. He was a man of
very agreeable conversation, and of great probity, as
well as immense erudition.—The following are the
titles of his principal works. 1. De claris interpres-
tibus, et de optimo genere interpretandi. 2. An edi-
tion of Origen's commentaries on the holy Scriptures,
in Greek and Latin. 3. A treatise on the origin
of the Romans. 4. Demonfrato evangelica, folio.
5. Quaestiones Alnetane de concordia rationis et fidei.
6. Of the situation of the terrestrial paradise, in
French. 7. A history of the commerce and navigati-
on of the ancients, which has been translated into
English. 8. Commentarius de rebus ad eum pertin-
entibus, \&c. 9. Hestiana. 10. Latin and Greek ver-

HUGHES (John), an ingenious and polite writer, was born in 1677. In the earliest parts of his youth, he cultivated the finer arts, poetry, drawing, and music, in each of which he by turns made a considerable progress; but followed those and his other studies only as agreeable amusements, under frequent confinement on account of his ill state of health. The lord Chancellor Cowper made him secretary for the commissions of the peace without his knowledge, and distinguished him with singular marks of his esteem. He continued in the same employment under the earl of Macclesfield, and held it to the day of his death; which happened in 1719, the very night in which his tragedy, entitled The Siege of Damascus, was first acted. He was then 42. He translated Fontenelle's dialogues of the dead, Vertot's revolutions of Portugal, and the letters of Abelard and Eloisa. He gave a very accurate edition of Spencer's works, with his life, a glossary, and remarks; and wrote several papers in the Tatler, Spectator, and Guardian. Mr Dunscombe, who married his sister, collected his poems and plays in 2 vols 12mo, in 1735.

\subsection*{HUGHILY (E0113)}
HUGHILY, a town of Afia, in the kingdom of
Bengal, seated on the most westerly branch of the ri-
ver Ganges. It is of large extent, reaching about
two miles along the river-side, and drives a great trade
in all the commodities of that country; affording rich
cargoes for 50 or 60 ships annually, besides what is
brought on carriages to the neighbouring towns.
Saltpetre is brought hither from Patna in vessels above
50 yards long and five broad. The inhabitants are
chiefly Indians; but there are also Portuguese, Eng-
lish, and other Europeans. E. Long 87. 55. N.
Lat. 22. 0.

HUGO Caper, chief of the third race of the kings
of France, being count of Paris and Orleans; he was
raised to the throne for his military valour and public
virtues in 987. See France, n 38.

\subsection*{HUGONIA (E0114)}
HUGONIA, in botany : A genus of the decandria
order, belonging to the monadelphia class of plants ;
and in the natural method ranking with those of which
the order is doubtful. The corolla is pentapetalous ;
the fruit is a plum with a triated kernel.

\subsection*{HUGUENOTS (E0115)}
HUGUENOTS, an appellation given by way of
contempt to the reformed or Protestant Calvinists of
France.

The name had its first rise in 1560; but authors are


\section*{Source record}
\begin{description}
\item[Scan ID] \texttt{britannica:309f765b1ecc83e0a7f7762e44b4a927a7555443:train:1327}
\item[Image file] \texttt{page-1327-7d41b1db21d5.jpg}
\item[Image SHA-256] \texttt{7d41b1db21d5e291f0cbf58aa92f8820d11602e700d5e337935c65b9d75b6211}
\item[Dataset] NationalLibraryOfScotland/encyclopaedia\_britannica\_illustrated at revision \texttt{309f765b1ecc83e0a7f7762e44b4a927a7555443} (CC0-1.0)
\item[Extraction] model \texttt{nvidia/nemotron-3.5-parse}, raw reply \texttt{extract/raw/page-1327-7d41b1db21d5.json}
\item[Raw reply SHA-256] \texttt{e9a8fec931cb4ddc0c6bbb98c0962f71a501f5bc68007dd0e2a1a5f9d0871550}
\item[Correction] model \texttt{nvidia/Nemotron-3.5-Super-VL}, block crops, raw replies \texttt{correct/raw/crop\_nothink/page-1327-7d41b1db21d5-b*.json}; parse: 1, super\_vl: 13
\item[Running header] none
\end{description}
\end{document}
